\section*{Chapitre \ref{ch:rc:operational-risk-and-rogue-trading} --- Risque opérationnel et fraude de trader (rogue trading)}

\begin{solution}{exo:rc:operational-risk-and-rogue-trading:1}
$0.12\times1+0.15\times9 = \text{EUR}~1.47$ milliard.
\end{solution}

\begin{solution}{exo:rc:operational-risk-and-rogue-trading:2}
Alors $(\mathrm{LC}/\mathrm{BIC})^{0.8}=1$ et $\mathrm{ILM}=\ln(e-1+1)=\ln e=1$~: une banque dont les pertes sont en ligne avec sa
taille conserve le BIC. Davantage de pertes relèvent le multiplicateur, moins de pertes le font baisser, les deux de façon moins que proportionnelle.
\end{solution}

\begin{solution}{exo:rc:operational-risk-and-rogue-trading:3}
Le règlement différé (dates de règlement très au-delà des deux jours de marché), associé à la règle de comptabilisation tardive, puisque la FSA
les a trouvés tous deux dans le dispositif.
\end{solution}

\begin{solution}{exo:rc:operational-risk-and-rogue-trading:4}
Un collatéral égal à un dixième de l'argent prêté signifie soit que les clients ont obtenu un crédit non garanti que personne
n'a approuvé, soit que l'argent n'allait pas du tout aux clients. Dans les deux cas, les demandes de financement qui ne peuvent
être réconciliées avec les positions et les clients connus doivent être suspendues jusqu'à explication.
\end{solution}

\begin{solution}{exo:rc:operational-risk-and-rogue-trading:5}
Les règles signalent des transactions authentiques pour des raisons sans rapport (une comptabilisation tardive ici, une exécution hors marché là),
si bien que deux règles coïncident rarement sur une transaction authentique~; un dispositif de dissimulation enfreint plusieurs règles à la fois par construction.
Sur le blotter, les fausses alertes passent de 7.1~\% à 0.05~\%, tandis que les détections ne passent que de 88.3~\% à 61.7~\%.
\end{solution}

\begin{solution}{exo:rc:operational-risk-and-rogue-trading:6}
Une position dissimulée exige que ses \omterm{def:rc:operational-risk-and-rogue-trading:rogue}{opérations fictives} soient reconduites, modifiées et expliquées chaque jour. Avec le trader
absent et tenu à l'écart, les positions sont gérées par quelqu'un d'autre, et les écarts, les appels de marge et les transactions non appariées
apparaissent.
\end{solution}

\begin{solution}{exo:rc:operational-risk-and-rogue-trading:7}
\texttt{audit\_demo()} renvoie \texttt{True} avant la modification et \texttt{False} après. Réécrire la chaîne à partir de l'entrée modifiée
produit une chaîne cohérente avec un hachage final différent~; si ce hachage a été copié ailleurs (dans un autre système, un rapport quotidien,
chez un tiers), la réécriture est détectée par comparaison.
\end{solution}

\begin{solution}{exo:rc:operational-risk-and-rogue-trading:8}
Des profits élevés et réguliers avec un risque mesuré faible sont la signature d'une position cachée ou d'une couverture fictive~:
les rendements sont trop beaux pour le risque que voit la société. La bonne réponse consiste à expliquer le P\&L par stratégie et par risque avant de relever la moindre limite.
\end{solution}

\begin{solution}{pb:rc:operational-risk-and-rogue-trading:1}
\textbf{1.} $6.4/49 = 13.1~\%$ (les 4.9 milliards nets intègrent les gains de 2007).
\textbf{2.} $6.4/0.05 = \text{EUR}~128$ milliards.
\textbf{3.} Vendre des dizaines de milliards d'actions en quelques jours représente une part importante du volume du marché~; les cours
se sont retournés contre la banque sous l'effet de ses propres ventes.
\textbf{4.} $49\times0.25/\sqrt{250}\times2.326 = \text{EUR}~1.80$ milliard.
\textbf{5.} Les \omterm{def:rc:operational-risk-and-rogue-trading:rogue}{opérations fictives} compensent les opérations réelles dans les registres de la société, de sorte que la position nette
vue par la VaR était faible.
\textbf{6.} Des transactions en sens inverse avec des contreparties internes ou fictives, ou des opérations à terme qui
ne seraient jamais réglées.
\textbf{7.} Des écarts internes non appariés, un règlement différé, une annulation et une modification à l'approche du reporting, et une comptabilisation tardive.
\textbf{8.} Rien, puisque les transactions étaient annulées avant leur date de confirmation~; une confirmation à la comptabilisation aurait fait apparaître des transactions qu'aucune contrepartie n'avait reconnues.
\textbf{9.} Marge et flux de trésorerie sur les positions réelles, bien supérieurs à la position nette expliquée.
\textbf{10.} Sans le trader pour reconduire et modifier les \omterm{def:rc:operational-risk-and-rogue-trading:rogue}{opérations fictives}, celles-ci se rompent.
\textbf{11.} Détection de 88.3~\% et fausses alertes de 7.1~\% avec une règle quelconque~; 61.7~\% et 0.05~\% avec deux règles.
\textbf{12.} Environ 200 par an avec une règle (142 fausses, 53 réelles), environ 40 avec deux, pour un blotter de
2\,060 transactions~; des fonctions de contrôle indépendantes du desk les instruisent.
\textbf{13.} Trop de fausses alertes, clôture sur la seule explication de la personne concernée sans preuve, et
absence d'escalade des récidives.
\textbf{14.} Il enregistre chaque comptabilisation, modification et clôture avec son horodatage et son auteur, et prouve que
le registre n'a pas été modifié.
\textbf{15.} Marge sur les positions, déclarations de positions aux bourses comparées au registre, et brut contre net.
\textbf{16.} Une première perte, dissimulée pour en éviter les conséquences, croît à mesure que le trader parie pour la récupérer.
\textbf{17.} Une règle sur les transactions annulées dont les dates de valeur sont éloignées, à la comptabilisation~: la confirmation arrivait trop tard.
\textbf{18.} Par la composante de pertes~: les pertes entrent dans la moyenne pendant dix ans et relèvent l'ILM.
\textbf{19.} Résultat à retenir~: \emph{l'arithmétique de janvier 2008}~: une perte de dénouement de 6.4 milliards d'euros sur
49 milliards est un mouvement défavorable moyen de 13.1~\%~; le premier contrôle de la chaîne qui l'aurait détecté se situe à la
comptabilisation~: signaler les transactions dont les dates de valeur sont éloignées et qui sont annulées avant confirmation.
\textbf{20.} Le suivi est assuré~: chaque alerte est clôturée sur preuves, indépendamment des personnes qu'elle contrôle.
\end{solution}

\begin{solution}{iq:rc:operational-risk-and-rogue-trading:1}
Le risque de perte résultant de processus, de personnes, de systèmes défaillants ou d'événements externes, y compris le risque juridique.
Les risques de marché et de crédit sont pris délibérément en contrepartie d'un rendement~; le \omterm{def:rc:operational-risk-and-rogue-trading:oprisk}{risque opérationnel} accompagne l'activité,
n'offre aucun gain possible et présente une queue épaisse.
\iqlookfor{la définition et la différence de nature.}
\end{solution}

\begin{solution}{iq:rc:operational-risk-and-rogue-trading:2}
\omterm{def:rc:operational-risk-and-rogue-trading:sod}{Séparation des tâches}, comptabilisation et confirmation indépendantes, réconciliation des liquidités, des marges et des
positions, explication du P\&L, limites, \omterm{def:rc:operational-risk-and-rogue-trading:sod}{congés obligatoires consécutifs} (block leave), et alertes instruites. Chez Barings, le trader dirigeait
le back-office, le financement n'était pas réconcilié avec les positions, et des profits exceptionnels n'étaient pas remis en question.
\iqlookfor{les contrôles et les défaillances documentées.}
\end{solution}

\begin{solution}{iq:rc:operational-risk-and-rogue-trading:3}
Flux issus de la comptabilisation, de la confirmation, du règlement-livraison et des données de marché~; règles comparables à celles du chapitre
avec, en plus, des profils de comportement par trader~; scores et seuils calibrés sur les fausses alertes~; gestion des dossiers avec preuves et
escalade~; un journal d'audit en ajout seul~; une propriété indépendante.
\iqlookfor{données, règles, scoring, instruction et indépendance.}
\end{solution}

\begin{solution}{iq:rc:operational-risk-and-rogue-trading:4}
Ne pas l'ignorer~; demander au responsable du desk et au risque l'explication par stratégie, et faire remonter aux fonctions de contrôle si aucune
n'arrive. Un profit régulier et inexpliqué est un signal d'alerte.
\iqlookfor{escalade, et non enquête menée de son propre chef.}
\end{solution}

\begin{solution}{iq:rc:operational-risk-and-rogue-trading:5}
Les fonds propres sont égaux au BIC (coefficients marginaux de l'indicateur d'activité) multiplié par l'ILM calculé sur dix ans de
pertes. Faiblesses~: la taille est un proxy grossier du risque, les pertes passées prédisent mal les événements de queue, et le dispositif
incite peu à améliorer les contrôles.
\iqlookfor{la formule et ses limites.}
\end{solution}

\begin{solution}{iq:rc:operational-risk-and-rogue-trading:6}
Trop de fausses alertes, clôture sur l'explication du trader, absence d'escalade des récidives, et enquêteurs dépendants du desk. Remède~: scores combinés,
preuves exigées pour clôturer, suivi des récidives, et
enquêteurs indépendants disposant d'une autorité reconnue.
\iqlookfor{fatigue d'alerte et indépendance.}
\end{solution}
